% ==================== 一、问题重述与分析 ====================
\section{问题重述与分析}

\subsection{问题背景}

微网是为解决光伏等分布式能源本地消纳、减少并网问题而发展起来的小型发用电系统。
出于经济性考虑，\textbf{非孤岛式微网}可在外部电网（下称外网）电价较低、或分布式能源有剩余电量时
为储能设备充电，并在外网电价较高时释放储存的电能，从而降低购电成本。
因此，如何在光伏出力不确定、电价时变的条件下制定\textbf{购电与充放电策略}，
是微网经济运行的核心问题。

本文研究的微网含光伏发电与储能设备，配置如下（题附录 1）：
最大容量 $12000\kw$，最大充放电功率 $5000$ kW，电量须保持在 $[1200,\ 10800]\kw$，
充放电单向效率均为 $90\%$，$2025$ 年 $1$ 月 $1$ 日 $0$:00 初始电量为 $6000\kw$。
微网不得向外网售电；微网供电不足时须向外网\textbf{紧急购电}，电价为交易时刻电价的 $5$ 倍。

\subsection{数据说明}

\begin{table}[H]
\centering
\caption{附件数据汇总}
\label{tab:data}
\small
\begin{tabularx}{\textwidth}{@{}c l X c@{}}
\toprule
附件 & 内容 & 用途 & 时间粒度 \\
\midrule
1 & 某天 144 段的电价、小区负载、光伏发电预测功率 & 问题一全部；二/三问的电价与预测回退 & 10 min \\
2 & 2025.1.1--12.31 每天的实际负载与光伏实际功率 & 二/三/四问的预测与结算 & 10 min \\
3 & 每日 $0$、$6$、$12$、$18$ 时发布的未来 24 h 整点光伏预报 & 三问、四之 4-3 的调整依据 & 1 h（整点）\\
4 & 2025.1.1--12.31 每天不同时间的电价 & 四问的结算电价 & 10 min \\
5 & 结果文件模板 & 各问 \texttt{result*.xlsx} 格式 & — \\
\bottomrule
\end{tabularx}
\end{table}

\subsection{统一口径约定}

为避免四问之间混淆，全文采用以下一致口径，后续各问不再重复声明。

\begin{enumerate}[itemsep=2pt,topsep=3pt]
  \item \textbf{时段划分}：一天划分为 $T=144$ 个时段，每段时长 $\Delta t=\tfrac{1}{6}$ h。
        附件中负载、光伏功率的单位为 kW，须乘以 $\Delta t$ 折算为每段电量（kWh）。
  \item \textbf{时间标签}：附件中的时刻为\textbf{区间末}时刻，即 \texttt{00:10} 表示 00:00--00:10，
        \texttt{0:00+1} 表示当天 24:00。
  \item \textbf{储能约束}：储电量 $E_t\in[1200,\ 10800]$；每段充、放电量
        $c_t,d_t\le 5000\Delta t=833.33\kw$；效率递推为
        $E_t=E_{t-1}+0.9c_t-d_t/0.9$。
  \item \textbf{交易规则}：不外送电；允许弃光与弃置多余的计划购电；缺口按 $5$ 倍电价紧急购电。
\end{enumerate}

\subsection{四问的递进关系}

四问在不确定性维度上层层递进，构成一条完整的建模主线：

\begin{center}
\begin{tabular}{@{}c c c c@{}}
\toprule
问题 & 电价 & 光伏信息 & 决策时点 \\
\midrule
问题一 & 已知（典型日）   & 已知              & 单时点（日内一次）\\
问题二 & 每日固定         & 仅历史预测        & 仅 $0$:00 \\
问题三 & 每日固定         & 历史预测 $+$ 发布预报 & $0/6/12/18$ 时 \\
问题四 & \textbf{实时波动} & 历史预测 $+$ 发布预报 & 4-2：仅 $0$:00；4-3：$0/6/12/18$ 时 \\
\bottomrule
\end{tabular}
\end{center}

相应地，模型从\textbf{确定性 LP}（问题一）过渡到\textbf{两阶段随机规划}（问题二），
再到\textbf{多阶段场景树随机规划}（问题三），最后引入\textbf{电价随机性}形成
双重不确定下的联合场景模型（问题四）。

\subsection{本文的主要特色}

\begin{enumerate}[itemsep=2pt,topsep=3pt]
  \item \textbf{报童结构的识别}：指出问题二的本质是报童问题，临界分位 $F^*=0.80$，
        为计划购电量提供了清晰的经济学解释，而非盲目调参。
  \item \textbf{LP 松弛紧性的严格证明}：证明问题一无需充放电互斥的二元约束，
        大幅降低计算复杂度（见 \S\ref{sec:q1-tight}）。
  \item \textbf{计费规则的统一化建模}：把问题三的分段计费规则写成单一线性表达式，
        并给出系数符号的必要性论证（见 \S\ref{sec:q3-fee}）。
  \item \textbf{场景树保证非前视}：节点内共享控制变量，使模型在结构上不可能"偷看"未来信息。
  \item \textbf{电价—净负载联合场景}：识别出期望不可分解性
        $\E[p\cdot(N-z)_+]\neq\E[p]\E[(N-z)_+]$，据此构造配对场景。
\end{enumerate}
